% !TEX root = CoeneDylanBP.tex
%%=============================================================================
%% Conclusie
%%=============================================================================

\chapter{Conclusie}%
\label{ch:conclusie}

Dit onderzoek had als doel een proof-of-concept te ontwikkelen van een kosteneffectief alternatief voor commerciële HDTS-systemen, op basis van computer vision. De centrale onderzoeksvraag luidde: \textit{Hoe kan een goedkoop computer vision-systeem worden ontwikkeld dat horizontale verplaatsing en vluchttijd bij trampolinespringen betrouwbaar kan benaderen, als betaalbaar alternatief voor het HDTS-meetsysteem?}

De ontwikkelde pipeline toont aan dat dit technisch haalbaar is met uitsluitend open-source modellen, een standaard videostream en een vaste cameraopstelling kunnen HD-landingszones automatisch worden onderscheiden en kan een vluchttijdschatting worden gemaakt. De resultaten en beperkingen van dit systeem worden hieronder per deelvraag samengevat.

\section*{Antwoord op de deelvragen}

\paragraph{Welke stappen vormen de voornaamste bottleneck?}
Zoals besproken in Hoofdstuk~\ref{ch:resultaten} en~\ref{ch:discussie} zijn twee structurele bottlenecks geïdentificeerd. De eerste is de framerate-afhankelijkheid: bij 25 fps bedraagt de onzekerheid per gedetecteerd landingsmoment 40\,ms, wat over tien sprongen cumulatief oploopt tot maximaal 800\,ms. Dit verklaart grotendeels de hoge ToF-afwijking (MAE = 5,253\,s). De tweede bottleneck is de opeenstapeling van fouten langs de diepte-as van de homografietransformatie: een kleine afwijking in de gedetecteerde hoekpunten resulteert door de perspectiefvervorming in een relatief grote metrische fout langs die as, wat de 2D-variant van de HD-schatting significant benadeelt (MAE = 1,806 vs. 0,649 in 1D-modus).

\paragraph{Welke computer vision-technieken zijn het meest geschikt?}
Op basis van de iteratieve vergelijking in Hoofdstuk~\ref{ch:methodologie} bleek een combinatie van drie technieken het meest geschikt. Een op YOLO11n-pose gebaseerde hoekpuntdetectie zorgt voor de geometrische kalibratie van het trampolineoppervlak. Een binaire classifier detecteert het landingsmoment, en SAM-2-segmentatie bepaalt de exacte contactpositie op het doek. Deze combinatie is robuuster dan de eerder onderzochte pose estimation-aanpak via ViTPose, die gevoeliger bleek voor occlusies en moeilijk te generaliseren was naar wisselende camerastandpunten. Optische stroming bleek geen meerwaarde te bieden voor de specifieke taak van landingsdetectie en werd na evaluatie niet opgenomen in de definitieve pipeline.

\paragraph{Welke pre-processing stappen zijn noodzakelijk?}
De pre-processing bestaat uit twee cruciale stappen, beschreven in Hoofdstuk~\ref{ch:methodologie}. Ten eerste is hoekpuntdetectie noodzakelijk om de homografietransformatie op te stellen: zonder nauwkeurige referentiepunten op het trampolineframe kan geen correcte omzetting van pixel- naar metrische coördinaten worden gemaakt. Ten tweede vereist de landingspositiebepaling via SAM-2 een uitsnede van het relevante beeldgebied, gestuurd door de bounding box van de YOLO-detector, om de segmentatie te richten op de springer en niet op andere personen in beeld.

\paragraph{Welke state-of-the-art modellen dragen bij aan hogere nauwkeurigheid?}
De vergelijking van modellen in Hoofdstuk~\ref{ch:methodologie} toont dat YOLO11n-pose, een lichtgewicht variante van het YOLO-framework, sterk presteert op hoekpuntdetectie met een detectieratio van 97,1\% en een gemiddelde genormaliseerde fout van 7,89\% over de testset. SAM-2 levert een nauwkeurigere segmentatie van de springer op dan een pure bounding-box benadering, en is robuuster bij varierende belichting en camerastandpunten. ViTPose werd als alternatief voor pose estimation geëvalueerd maar leverde in de praktijk geen betere resultaten op dan de YOLO-gebaseerde aanpak, mede door de hogere gevoeligheid voor occlusie in wedstrijdomstandigheden.

\paragraph{Is gescheiden detectie effectiever dan een end-to-end model?}
De modulaire pipeline, waarbij landingsmomentdetectie en landingspositiebepaling als afzonderlijke stappen zijn uitgewerkt, biedt als voordeel dat elke component afzonderlijk kan worden geëvalueerd en verbeterd. De resultaten in Hoofdstuk~\ref{ch:resultaten} bevestigen echter ook de keerzijde: fouten in vroege stappen werken door naar latere stappen, en het is lastig te isoleren welke component verantwoordelijk is voor de uiteindelijke afwijking. Een end-to-end model dat beide taken gelijktijdig oplost, zou deze foutdoorwerking kunnen beperken, maar vereist gelabelde per-sprong grondwaarheid die in de gebruikte dataset ontbrak. De vergelijking blijft daardoor theoretisch.

\section*{Antwoord op de hoofdvraag}

De hoofdvraag kan bevestigend worden beantwoord: een goedkoop computer vision-systeem \textit{kan} horizontale verplaatsing en vluchttijd bij trampolinespringen benaderen. De proof-of-concept pipeline onderscheidt HD-landingszones betrouwbaar in 1D-modus (MAE = 0,649) en maakt vluchttijdschatting mogelijk op basis van frame-telling. De absolute nauwkeurigheid is echter nog onvoldoende om het HDTS-systeem te vervangen. De ToF-MAE van 5,253\,s overschrijdt de gestelde eis van 0,2\,s aanzienlijk, en ook de HD-scores wijken in absolute zin te veel af voor directe inzet in competitieverband.

Toch toont dit onderzoek aan dat de fundamentele bouwstenen beschikbaar en haalbaar zijn. De pipeline draait volledig op open-source modellen en lichtgewicht hardware, wat de kostprijs potentieel ruimschoots onder die van commerciële systemen brengt. De voornaamste beperkingen zijn niet inherent aan de gekozen technologieën, maar aan de beschikbare testomstandigheden. Het ontbreken van per-sprong grondwaarheid en de lage framerate van de gebruikte wedstrijdvideo's beperken zowel de nauwkeurigheid als de evaluatiemogelijkheden. Met een hogere framerate, een tweede camera voor de diepte-as en een dataset met HDTS-annotaties zou een significant nauwkeuriger systeem realiseerbaar zijn, zoals beschreven in Hoofdstuk~\ref{ch:discussie}.

De bijdrage van dit onderzoek aan het vakgebied ligt in het aantonen dat een iteratief opgebouwde, modulaire computer vision-pipeline een geloofwaardig startpunt vormt voor een betaalbaar meetsysteem in trampolinespringen. Daarmee biedt het een concreet alternatief perspectief voor turnclubs die zich de kostprijs van het HDTS-systeem niet kunnen veroorloven.
